% Actual scalar refinement and stability; action/frame equations are separate.
\statement{A.68}{Proved}{The full cubic scalar kernel has dimension 128 and its next refinement is stable}
At a finite integral level n of the actual quotient in A.66, let
\(W_n\) be its source kernel and put
\(W_{n,1}=[W_n,\Pi]\), with closed commutators understood.
Then
\[
 W_n/W_{n,1}\simeq C_2^{128},\qquad
 [W_{n,1},\Pi]=W_{n,1}.
 \tag{A.68.1}
\]
The scalar identifications are compatible with the integral transition
maps. They also hold after adjoining any earlier finite retained family.
If \(W_\infty=\bigcap_nW_n\), then
\[
 W_\infty/[W_\infty,\Pi]\simeq C_2^{128},\qquad
 [[W_\infty,\Pi],\Pi]=[W_\infty,\Pi].
 \tag{A.68.2}
\]
These statements identify the actual scalar refinement. The action,
negative square, and normalizer equations on that refinement require
separate verification.

\begin{proof}
\emph{Reduction to actual finite central extensions.}\par
Write \(Q_n\) for the finite residual over the initial C at the
chosen level. Its abelianization is
\(A_n=L_n^4\oplus V^4\), and its derived group has blocks
\[
 M_n,\qquad T=(V^{(2)})^6,\qquad
 U=(V^{(3)})^4,\qquad Z=(V^{(3)})^{16}.
 \tag{A.68.3}
\]
Here \(M_n\) is the old globally central group of A.58--A.62.
Only T has nonzero commutators with the abelianization; their
images are U. The group has class three. Let N be the actual
normal odd weight subgroup in the finite head H. We have
\((A_n)_N=0\).

Consider any finite image of the first scalar refinement that retains
\(Q_n\) as the quotient of its R-image. Such images are cofinal for
detecting the scalar kernel: refine any finite quotient by a finite
quotient of \(\Pi/W_n\) faithful on its finite normal subgroup
\(Q_n\). Denote the resulting R-image by \(\widetilde Q\),
and its new globally central kernel over \(Q_n\) by B.
The original-base vanishing \(H^1(C,A_n^\vee)=0\), together
with \(R=[R,\Pi]\), gives
\[
 \widetilde Q^{\rm ab}=A_n.
 \tag{A.68.4}
\]
Indeed, the dual of a new abelianized central kernel is globally
fixed, whereas the dual of \(\widetilde Q^{\rm ab}\) has no
fixed vectors. The invariant exact sequence injects that new dual
into \(H^1(C,A_n^\vee)\). This is the actual Pontryagin-dual
argument of A.53. An odd-primary part of B splits as a central
factor over the finite two-group \(Q_n\), and would survive
abelianization. Hence \(\widetilde Q\) is a finite two-group.

\emph{The fourth lower-central subgroup vanishes.}\par
The entire inverse image of \(M_n\) centralizes
\(\widetilde Q\). Its commutators lie in B and define an
equivariant pairing
\(M_n\otimes A_n\longrightarrow B\). The first module
is globally trivial, and \((A_n)_N=0\), so the pairing is zero.
The same observation applies to any old globally central scalar
block in a subsequent refinement.

Since \(Q_n\) has class three,
\(\gamma_4(\widetilde Q)\subset B\) and
\(\gamma_5(\widetilde Q)=1\). Its possible generators are
commutators of lifts of U with \(A_n\). These pairings factor
through binary modules. Extend scalars to \(k=\mathbb F_8\)
and decompose by the actual N-weights. A U-weight \(3\lambda\)
can have a scalar partner only at weight
\(-3\lambda=4\lambda\).

Represent this U-weight by
\([[a_i(\lambda),a_j(\lambda)],a_k(\lambda)]\).
Put \(t=[a_i,a_j]\), of weight \(2\lambda\), and let d
have weight \(4\lambda\). The bracket \([t,d]\) projects
to zero: the old triple bracket requires the same Frobenius
family. Its lift is scalar of nontrivial weight \(6\lambda\),
so it is zero. The class-four Hall--Witt identity reduces
\([
 [[a_i,a_j],a_k],d
]\) to \([t,[a_k,d]]\), up to the irrelevant sign convention.
The inner bracket projects to the ordinary Z-block, which is
central in \(Q_n\). Its lifts have central commutators with
\(\widetilde Q\), and therefore centralize the derived subgroup.
Thus \([t,[a_k,d]]=1\). This kills \(\gamma_4\).

Furthermore \(\gamma_3(\widetilde Q)\) is elementary abelian.
If g is derived, its square projects into \(M_n\), since the
other three blocks in (A.68.3) have exponent two. The full inverse
image of \(M_n\) is central, and hence
\[
 [g,x]^2=[g^2,x]=1.
 \tag{A.68.5}
\]
This proves the exponent assertion for the generators of
\(\gamma_3\).

\emph{No additional scalar survives in degree two.}\par
Pass to \(\widetilde Q/\gamma_3(\widetilde Q)\).
Its derived group D is a quotient of
\(\bigwedge^2_{\mathbb Z}A_n\), and surjects onto
\(M_n\oplus T\oplus Z\). Coprime averaging over N makes
\(D^N\) surject onto \(M_n\); its kernel is exactly the image
of B in this class-two quotient. Both subquotients are globally
trivial. Automorphisms fixing them form a subgroup of their Hom
group, so the head action on \(D^N\) has two-group image.
The action factors through the actual H-action on \(A_n\),
and thus through \(H/O^2(H)\). Consequently \(D^N\) is a
quotient of the already exact module
\[
 \left(\bigwedge^2_{\mathbb Z}A_n\right)_{O^2(H)}=M_n.
 \tag{A.68.6}
\]
It also surjects onto \(M_n\). Equality of finite orders forces
its new scalar kernel to vanish. Therefore
\(B\subset\gamma_3(\widetilde Q)\), so B is elementary
abelian and is generated by actual third commutators.

\emph{An upper bound from actual triple pairings.}\par
A first bracket in \(M_n\) gives zero. A first bracket in T
has weight \(2\lambda\) from two equal-family inputs at
\(\lambda\); its only scalar third partner has weight
\(-2\lambda\). This is the opposite physical weight, so both
alternative brackets in Jacobi project to zero. Their lifts are
central. The T contribution is therefore zero.

An ordinary Z bracket uses free-row inputs at weights
\(\lambda,2\lambda\) at the same physical weight. Its scalar
third partner has weight \(4\lambda\). If that partner is a
binary s-row, both alternative brackets vanish and Jacobi again
kills the result. The only remaining triples use three free rows
in the three distinct Frobenius families
\(\lambda,2\lambda,4\lambda\).

For one physical weight there are three cyclic bracket tensors
on \(k^4\otimes k^4\otimes k^4\), with one independent Jacobi
relation for each ordered row triple. Their dimension is at most
\[
 3\cdot4^3-4^3=128
 \quad\hbox{over }k.
 \tag{A.68.7}
\]
The subgroup \(O^2(H)/N\), of order 36, acts regularly on the
physical weights and identifies these tensor spaces. The scalar
product has weight seven, hence is one. The remaining order-two
stabilizer acts trivially: it fixes the chosen weight line with
scalar one and does not mix the cusp multiplicities. Descent from
k preserves dimension. Thus (A.68.7) is a bound of 128 over
\(\mathbb F_2\), not 384. This argument uses actual commutator
pairings and Hall--Witt in a finite class-three group, not a
presumed realization of a graded Lie algebra.

\emph{A cover attaining all 128 directions.}\par
On the ordinary class-two quotient write
\[
 (a,z)(b,w)=
 \left(a+b,z+w+(a_i^2b_j)_{i,j}\right),
 \qquad a,b\in k^4,\quad z,w\in k^{16}.
\]
For each ordered triple define
\[
 F_{ijk}((a,z),(b,w))=z_{ij}b_k^4+a_i^2w_{jk}^4.
 \tag{A.68.8}
\]
These are k-valued central cocycles: the two coboundary terms
cancel because \((b_j^2c_k)^4=b_jc_k^4\). All 64 together
have raw central group \(k^{64}\), of binary dimension 192.
Define
\[
 (\sigma X)_{ijk}=X_{jki}^4.
 \tag{A.68.9}
\]
Then \(\sigma^3=1\), and a raw triple commutator has value
\[
 (1+\sigma)P,\qquad
 P_{ijk}=(a_i^2b_j+b_i^2a_j)c_k^4.
 \tag{A.68.10}
\]
Ordinary commutators span all \(k^{16}\); multiplying by arbitrary
third inputs spans all \(k^{64}\). Hence the raw third subgroup
is exactly \(\operatorname{im}(1+\sigma)\).

The fixed space of \(\sigma\) has binary dimension 64. Its four
diagonal index triples each contribute one fixed binary coordinate;
the remaining 60 triples form 20 cyclic orbits, each contributing
one arbitrary k coordinate. Since the order of \(\sigma\) is odd,
its fixed space complements \(\operatorname{im}(1+\sigma)\).
Quotient by that fixed space, equivalently project by the idempotent
\(\sigma+\sigma^2\). This produces a stem central cover of
binary dimension 128. Independent coordinatewise quotients
\(k/\mathbb F_2\) would be incorrect for the mixed triples:
they lose 40 directions. The raw exact commutator certificate
checks dimension 128 and the stated fixed-space complement:
\path{work_CE/code/check_generalized_scalar_triples.py} and
\path{work_CE/generalized_scalar_triple_certificate.json}.

Take the 36 physical weight copies and push out their central
factors by summation. The degrees seven and fourteen in (A.68.8)
make this H-equivariant. Pair the result with the actual \(Q_n\).
The new scalar directions are attained by third commutators.
Coprime averaging on the elementary third subgroup isolates their
N-fixed component from the old U, which has no N-fixed vectors.
Thus the complete new central kernel lies in the Frattini subgroup.
The source still surjects onto the complete original eight-cocycle
Frattini quotient. Finite Frattini therefore proves that the
fiber product is an actual source quotient. Its attained scalar
kernel has dimension 128. The upper bound proves the first
assertion of (A.68.1).

\emph{The next scalar refinement is stable.}\par
In a further scalar refinement, the entire inverse image of the
old scalar \(C_2^{128}\) centralizes the residual, by the same
pairing with \((A_n)_N=0\). In particular the ordinary Z-lifts
still lie in the second upper center: their commutators project
into that old scalar block, whose full preimage is now central.
They still centralize the derived subgroup. The fourth-commutator
argument and the degree-two exhaustion are unchanged.

The N-fixed part of the new elementary third subgroup is an
extension of the globally trivial old \(C_2^{128}\) by the
new globally trivial scalar kernel. Its head action has two-group
image and factors through the actual action on the triple pairings.
Thus it factors through \(H/O^2(H)\). The regular
\(O^2(H)/N\)-orbit identification and the trivial remaining
stabilizer give the same 128-dimensional upper bound. This group
already surjects onto the old 128-dimensional scalar group.
The new kernel must be zero. Cofinality of the finite central
quotients gives the second assertion of (A.68.1).

\emph{Integral limits and older retained families.}\par
The explicit scalar cover uses only the modulo-two cusp data.
Its central generators are therefore the same at every integral
level. The natural maps between the scalar kernels are surjective
and identify these 128 generators, hence are isomorphisms.
Set \(U_\infty=\bigcap_nW_{n,1}\). The finite-image
stabilization and compactness argument of A.62 proves
\(U_\infty=[U_\infty,\Pi]\). Compactness of the fibers over
each of the common 128 scalar values gives
\(W_\infty/U_\infty=C_2^{128}\). Consequently
\[
 U_\infty=[U_\infty,\Pi]
 \subset[W_\infty,\Pi]\subset U_\infty,
\]
which proves (A.68.2).

For an earlier finite retained family use its actual base
\(C_{\rm old}\). The established type separation and continuous
five-term sequence give \(H^1(C_{\rm old},A_n^\vee)=0\),
and the actual new residual remains \(Q_n\). All subsequent
averaging uses the unchanged finite H-action; it assumes no
new odd subgroup inside \(C_{\rm old}\). The preceding proof
therefore applies at every finite common refinement, and
compactness retains those families simultaneously. Nothing in
this scalar calculation supplies a solution of the separate
normalizer equations on the new cover.
\end{proof}
